\section{Chapter 1: General Vector Spaces}

\subsubsection*{Vector Addition Axioms}
\begin{enumerate}
    \item \textbf{Closure under Addition:} 
    For all $u, v \in V$, $u + v \in V$.
    
    \item \textbf{Commutativity:} 
    $u + v = v + u$ for all $u, v \in V$.
    
    \item \textbf{Associativity:} 
    $(u + v) + w = u + (v + w)$ for all $u, v, w \in V$.
    
    \item \textbf{Additive Identity:} 
    There exists $0 \in V$ such that $v + 0 = v$ for all $v \in V$.
    
    \item \textbf{Additive Inverse:} 
    For each $v \in V$, there exists $-v \in V$ such that $v + (-v) = 0$.
\end{enumerate}

\subsubsection*{Scalar Multiplication Axioms}
\begin{enumerate}
    \setcounter{enumi}{5}
    \item \textbf{Closure under Scalar Multiplication:} 
    For all $c \in \mathbb{R}$ and $v \in V$, $cv \in V$.
    
    \item \textbf{Vector Distributivity:} 
    $c(u + v) = cu + cv$ for all $c \in \mathbb{R}$ and $u, v \in V$.
    
    \item \textbf{Scalar Distributivity:} 
    $(c + d)v = cv + dv$ for all $c, d \in \mathbb{R}$ and $v \in V$.
    
    \item \textbf{Associativity:} 
    $(cd)v = c(dv)$ for all $c, d \in \mathbb{R}$ and $v \in V$.
    
    \item \textbf{Multiplicative Identity:} 
    $1v = v$ for all $v \in V$.
\end{enumerate}

\subsection{The Elements of a Vector Space}
The elements of a vector space $V$ are called \textbf{vectors}. The elements of the field $\mathbb{R}$ are called \textbf{scalars}.

\vspace{1em}
\noindent\textbf{Examples of Vector Spaces:}

(1) $\mathbb{R}^n = \left\{ \begin{bmatrix} x_1 \\ \vdots \\ x_n \end{bmatrix} : x_1, \dots, x_n \in \mathbb{R} \right\}$ under standard vector addition and scalar multiplication.

\vspace{1em}

(2) $M_{m \times n}(\mathbb{R})$ under standard matrix addition and scalar multiplication.

\vspace{1em}

(3) $P_n(\mathbb{R})$, the space of all polynomials of degree $\le n$, defined with operations:
\[
(f + g)(x) = f(x) + g(x) \quad \text{and} \quad (cf)(x) = c f(x) \quad (c \in \mathbb{R}).
\]
Then $P_n(\mathbb{R})$ with these operations is a vector space over $\mathbb{R}$.

% =========================================================================
% SECTION 1.2: AXIOMATIC PROPERTIES & SUBSPACES
% =========================================================================

\subsection{Fundamental Properties Derived from the Axioms}

The definition of a vector space guarantees that standard arithmetic rules apply without relying on the specific nature of the objects (whether they are $n$-tuples, matrices, or functions).

\vspace{1em}

\noindent \textbf{Theorem (Elementary Arithmetic Properties of Vector Spaces):} Let $V$ be a vector space over $\mathbb{R}$. For all $\vec{x} \in V$ and all $t \in \mathbb{R}$:
\begin{enumerate}
    \item $0\vec{x} = \vec{0}$
    \item $t\vec{0} = \vec{0}$
    \item $(-1)\vec{x} = -\vec{x}$
    \item If $t\vec{x} = \vec{0}$, then either $t = 0$ or $\vec{x} = \vec{0}$.
\end{enumerate}

\bigskip

\begin{proof}
\textbf{Proof of (1):} In the real numbers, $0 + 0 = 0$. By scalar distributivity:
\[
0\vec{x} = (0 + 0)\vec{x} = 0\vec{x} + 0\vec{x}.
\]
Since $0\vec{x} \in V$, its additive inverse $-(0\vec{x})$ exists in $V$. Adding this inverse to both sides:
\begin{align*}
0\vec{x} + (-(0\vec{x})) &= (0\vec{x} + 0\vec{x}) + (-(0\vec{x})) \\
\vec{0} &= 0\vec{x} + \bigl(0\vec{x} + (-(0\vec{x}))\bigr) \\
\vec{0} &= 0\vec{x} + \vec{0} \\
\vec{0} &= 0\vec{x}.
\end{align*}
Therefore, $0\vec{x} = \vec{0}$.

\bigskip
\vspace{0.5em}

\textbf{Proof of (2):} By the additive identity axiom, $\vec{0} + \vec{0} = \vec{0}$. Multiplying by scalar $t$ and using \indent vector distributivity:
\[
t\vec{0} = t(\vec{0} + \vec{0}) = t\vec{0} + t\vec{0}.
\]
\indent Adding $-(t\vec{0})$ to both sides and associating terms yields $t\vec{0} = \vec{0}$.

\bigskip
\vspace{0.5em}

\textbf{Proof of (3):} We show that $(-1)\vec{x}$ behaves as the unique additive inverse of $\vec{x}$:
\[
\vec{x} + (-1)\vec{x} = 1\vec{x} + (-1)\vec{x} = (1 + (-1))\vec{x} = 0\vec{x} = \vec{0}.
\]
\indent By uniqueness of the additive inverse in $V$, $(-1)\vec{x} = -\vec{x}$.

\bigskip
\vspace{0.5em}

\textbf{Proof of (4):} Suppose $t\vec{x} = \vec{0}$. If $t = 0$, the claim holds. If $t \neq 0$, the multiplicative inverse \indent $t^{-1} \in \mathbb{R}$ exists. Multiplying both sides:
\[
t^{-1}(t\vec{x}) = t^{-1}\vec{0} \implies (t^{-1}t)\vec{x} = \vec{0} \implies 1\vec{x} = \vec{0} \implies \vec{x} = \vec{0}.
\]
\end{proof}

\subsection{Subspaces and the Subspace Test}

\noindent\textbf{Definition (Subspace).}
Let $V$ be a vector space over $\mathbb{R}$. A subset $U \subseteq V$ is called a \textbf{subspace} of $V$ if $U$ is itself a vector space over $\mathbb{R}$ under the addition and scalar multiplication operations inherited from $V$.

\medskip
\noindent\textbf{Theorem (Subspace Test).} 
Let $V$ be a vector space over $\mathbb{R}$. A subset $U \subseteq V$ is a subspace of $V$ if and only if:
\begin{enumerate}[itemsep=0.8em, topsep=0.5em]
    \item $\vec{0} \in U$, where $\vec{0}$ is the zero vector of $V$ (in particular, $U \neq \emptyset$).
    \item For all $\vec{x}, \vec{y} \in U$, $\vec{x} + \vec{y} \in U$ (\textbf{Closure under vector addition}).
    \item For all $\alpha \in \mathbb{R}$ and $\vec{x} \in U$, $\alpha \vec{x} \in U$ (\textbf{Closure under scalar multiplication}).
\end{enumerate}

\begin{proof}
\noindent\textbf{($\implies$)} If $U$ is a subspace, it satisfies all vector space axioms on its own. In particular, it must be closed under addition and scalar multiplication, and contain an additive identity $\vec{0}_U \in U$. 

Since $\vec{0}_U + \vec{0}_U = \vec{0}_U$, adding the inverse $-\vec{0}_U$ computed in $V$ gives $\vec{0}_U = \vec{0}_V$.

\bigskip
\vspace{0.5em}

\noindent\textbf{($\impliedby$)} Assume conditions (1), (2), and (3) hold. Closure under addition and scalar multiplication, as well as the existence of the additive identity, are directly satisfied. 

For any $\vec{x} \in U$, setting $\alpha = -1$ ensures that $-\vec{x} = (-1)\vec{x} \in U$.

The remaining algebraic axioms (commutativity, associativity, and distributivity) hold for all elements in $U$ automatically because they hold throughout the larger space $V$. Thus, $U$ is a vector space.
\end{proof}

\subsection*{Worked Examples: Subspace Analysis}

\noindent\textbf{Example 1 (Polynomial with Root Constraint).}
Let $V = P_2(\mathbb{R})$. Define the subset:
\[ 
U = \{p(x) \in P_2(\mathbb{R}) \mid p(3) = 0\}.
\]
Prove that $U$ is a subspace of $P_2(\mathbb{R})$.

\begin{proof}
We test the three conditions of the Subspace Test:
\begin{enumerate}[itemsep=0.8em, topsep=0.5em]
    \item \textbf{Zero Vector:} The zero polynomial is $z(x) = 0 + 0x + 0x^2$. Evaluating at $x = 3$ gives $z(3) = 0$. Hence, $z \in U$.
    \item \textbf{Closure under Addition:} Let $p(x), q(x) \in U$. Then $p(3) = 0$ and $q(3) = 0$. Consider their sum evaluated at $x = 3$:
    \[
    (p + q)(3) = p(3) + q(3) = 0 + 0 = 0.
    \]
    Thus, $p + q \in U$.
    \item \textbf{Closure under Scalar Multiplication:} Let $p(x) \in U$ and $\alpha \in \mathbb{R}$. Then:
    \[
    (\alpha p)(3) = \alpha \cdot p(3) = \alpha(0) = 0.
    \]
    Thus, $\alpha p \in U$.
\end{enumerate}
By the Subspace Test, $U$ is a subspace of $P_2(\mathbb{R})$.
\end{proof}

\medskip
\noindent\textbf{Example 2 (Traceless Matrices).} 
Let $V = M_{2\times 2}(\mathbb{R})$. Define:
\[
W = \{ A \in M_{2\times 2}(\mathbb{R}) \mid \operatorname{tr}(A) = 0 \}.
\]
Recall that for $A = \begin{bmatrix} a & b \\ c & d \end{bmatrix}$, $\operatorname{tr}(A) = a + d$. Show that $W$ is a subspace of $M_{2\times 2}(\mathbb{R})$.

\begin{proof}
\begin{enumerate}[itemsep=0.8em, topsep=0.5em]
    \item \textbf{Zero Vector:} The zero matrix $O_{2\times 2} = \begin{bmatrix} 0 & 0 \\ 0 & 0 \end{bmatrix}$ satisfies $\operatorname{tr}(O_{2\times 2}) = 0 + 0 = 0$. Hence, $O_{2\times 2} \in W$.
    \item \textbf{Closure under Addition:} Let $A, B \in W$, so $\operatorname{tr}(A) = 0$ and $\operatorname{tr}(B) = 0$. By trace linearity:
    \[
    \operatorname{tr}(A + B) = \operatorname{tr}(A) + \operatorname{tr}(B) = 0 + 0 = 0.
    \]
    Hence, $A + B \in W$.
    \item \textbf{Closure under Scalar Multiplication:} Let $A \in W$ and $\alpha \in \mathbb{R}$. Then:
    \[
    \operatorname{tr}(\alpha A) = \alpha \operatorname{tr}(A) = \alpha(0) = 0.
    \]
    Hence, $\alpha A \in W$.
\end{enumerate}
Thus, $W$ is a subspace of $M_{2\times 2}(\mathbb{R})$.
\end{proof}

\medskip
\noindent\textbf{Example 3 (Affine Coefficient Failure).} 
Let $V = P_2(\mathbb{R})$ and consider the subset:
\[
S = \{ p(x) = ax^2 + bx + c \in P_2(\mathbb{R}) \mid a - 2b + c = 3 \}.
\]
Determine whether $S$ is a subspace of $V$.

\medskip
\noindent\textbf{Solution:} 
The zero polynomial $z(x) = 0x^2 + 0x + 0$ has coefficients $a = 0, b = 0, c = 0$. Evaluating the defining equation:
\[
a - 2b + c = 0 - 2(0) + 0 = 0 \neq 3.
\]
Because the zero polynomial is not in $S$, condition (1) of the Subspace Test fails immediately. Thus, $S$ is \textbf{not} a subspace of $P_2(\mathbb{R})$.

% =========================================================================
% SECTION 1.3: SPAN, INDEPENDENCE, AND BASES
% =========================================================================

\subsection{Linear Span and Minimal Subspaces}